% Guard: macros.tex is \input both by main.tex (after the packages it needs,
% e.g. amsmath for \DeclareMathOperator) and by individual section files so
% they remain compilable on their own. Loading it twice would redefine every
% \newcommand and abort, so the second and later loads are no-ops.
\ifdefined\PQRmacrosloaded\endinput\fi
\let\PQRmacrosloaded\relax

% --- 0. revision marking ---
% Text added to Section 2 during the current revision is wrapped in \added{...}
% so the advisors can see at a glance what is new. Redefine to \relax to print
% the paper without the marking.
\newcommand{\added}[1]{\textcolor{blue}{#1}}
\long\def\added#1{\begingroup\color{blue}#1\endgroup}  % \long: may span paragraphs


% --- A. operator-style score and function names ---
\newcommand{\HD}{\mathit{HD}}                        % histogram distance; use \HD_\gamma
\newcommand{\Attr}{\mathit{Attr}}                    % attribute set of a schema
\newcommand{\Attrs}{\mathit{Attrs}}                  % attribute set argument
% Abbreviated in MATH only; the prose still spells the names out in full.
% Spelling them out inside formulas makes displayed equations far too wide
% for a two-column layout (the product-rule step below overflowed by 200pt).
\newcommand{\Use}{\mathit{Use}}
\newcommand{\Cov}{\mathit{Cov}}
\newcommand{\Rep}{\mathit{Rep}}
\newcommand{\DivS}{\mathit{Div}}
\newcommand{\TVD}{\mathit{TVD}}
\newcommand{\GS}{\mathit{GS}}                        % global sensitivity
\DeclareMathOperator{\Dom}{Dom}
\DeclareMathOperator{\Lap}{Laplace}

% --- B. core objects ---
\newcommand{\Qone}{Q_1}
\newcommand{\Qtwo}{Q_2}
\newcommand{\pD}{p(D)}
\newcommand{\views}{\mathcal{P}}
\newcommand{\hist}[2]{H_{#1,#2}}                     % \hist{\Qone(\pD)}{A}
\newcommand{\dhist}{\Delta H}
\newcommand{\nhist}{\widetilde{\Delta H}}            % the released noisy object

% --- C. privacy budget and parameters ---
\newcommand{\eps}{\epsilon}
\newcommand{\epsP}{\epsilon_P}
\newcommand{\epsS}{\epsilon_S}
\newcommand{\epsH}{\epsilon_H}
\newcommand{\gammaref}{\gamma_{\mathrm{ref}}}
\newcommand{\gammaeval}{\gamma_{\mathrm{eval}}}
\newcommand{\taumax}{\tau_{\max}}

% --- D. mechanisms ---
\newcommand{\ExpMech}{\textsc{ExpMech}}
\newcommand{\OneShotTop}[1]{\textsc{OneShotTop-}#1}
\newcommand{\AdaptiveTau}{\textsc{AdaptiveTau}}
% --- E. experiment baselines (used in sections/experiments.tex) ---
\newcommand{\Bone}{\ensuremath{\mathsf{B1}}}
\newcommand{\Btwo}{\ensuremath{\mathsf{B2}}}
\newcommand{\Bthree}{\ensuremath{\mathsf{B3}}}

\section{Algorithm}\label{sec:solution}

\subsection{Overview}
\label{subsec:algorithm-overview}

The algorithm first privately selects a small set of useful predicate views, then constructs explanation candidates by pairing selected views with attributes, and finally privately selects and releases the top-$k$ histogram explanations.
This structure mirrors recent differentially private explanation frameworks, which separate the definition of explanation quality from the design of low-sensitivity private selection mechanisms \cite{TaoGMR22}.



The measures the algorithm optimises, and the reasons they take the form they do, are those of Section~\ref{sec:explanation-quality}; their sensitivity is Proposition~\ref{prop:sensitivity-itemized}, and the two stages and their utilities are Section~\ref{subsec:two-stages}.
What remains here is the algorithmic part: how the two per-round parameters are derived from the privacy budget, why the greedy selection is exact, the two algorithms themselves, and the privacy and complexity analyses.



\subsection{Budget-Derived $\tau$ and $\gamma$}
\label{subsubsec:adaptive-params}

Both $\tau$ and $\gamma$ admit natural data-dependent choices, but selecting either privately would consume budget that the selection and release stages need.
We instead follow the budget-annealing strategy of AIM~\cite{McKennaMSM22}, in which hyper-parameters are functions of the \emph{privacy accountant} rather than of the data:
\begin{equation}
\tau_i=\tau_{\min}+\tfrac{\epsilon_{\text{used}}}{\epsilon_S}\left(\tau_{\max}-\tau_{\min}\right),
\qquad
\gamma_i=\gamma_{\text{ref}}\cdot\frac{\epsilon_{\text{ref}}}{\epsilon_i}.
\label{eq:adaptive-params}
\end{equation}
Both depend only on public quantities, so by post-processing (Theorem~\ref{thm:post-processing}) they incur no privacy cost.
The form of $\gamma_i$ follows from Proposition~\ref{prop:hd-sensitivity}: the exponential mechanism's ability to separate candidates scales as $\epsilon_i/(2\,\GS)=\epsilon_i(1+\gamma)/2$, so holding separation roughly constant as $\epsilon_i$ varies requires $\gamma_i\propto 1/\epsilon_i$ --- smooth harder when a round is noisy, less when it is sharp.

The reference value $\gamma_{\text{ref}}$ must still be fixed in advance from public information.
Section~\ref{subsec:gamma-tradeoff} shows that anchoring it to a fraction of $|D|$ is defensible but crude, because what $\gamma$ competes against is the size of individual histogram \emph{bins}, whose distribution relative to $|D|$ varies sharply across datasets.

\subsection{Exact Non-Private Selection}
\label{subsec:non-private-hardness}

Selection problems of this shape are commonly NP-hard, and an earlier version of this work asserted as much.
That is not the case for the objective the explanation stage optimises: Definition~\ref{def:setobj} separates over attribute groups, and with its scores held fixed the resulting problem is solvable exactly in polynomial time.

\begin{theorem}[Greedy is exactly optimal]
\label{thm:non-private-hardness}
Fix a candidate set $\mathcal{C}'\subseteq\mathcal{C}$, a smoothing constant $\gamma$, and a non-increasing $w$.
Selecting $E\subseteq\mathcal{C}'$ with $|E|=k$ maximising $F_w(E)$ is solvable in polynomial time, and the greedy procedure --- repeatedly adding the candidate of largest marginal gain --- returns an optimal solution.
\end{theorem}

\begin{proof}
By Definition~\ref{def:setobj}, $F_w(E)=\sum_A g_A(E)$ where $g_A$ depends only on the members of $E$ whose attribute is $A$.
There is therefore no interaction between groups, and a solution is determined by how many candidates $g_A\ge0$ it takes from each group, subject to $\sum_A g_A=k$, together with which ones.
For a fixed $g_A$, the best choice within group $A$ is its $g_A$ highest-scoring candidates, since $w$ is non-increasing and the scores are ranked (rearrangement inequality).
Hence the value of taking $g$ candidates from group $A$ is
\[
f_A(g)=\sum_{j=1}^{g} w_j\, h^A_{(j)},
\]
whose successive marginal gains are $f_A(g)-f_A(g-1)=w_g\,h^A_{(g)}$.
These are non-increasing in $g$, because $w_g$ is non-increasing by assumption and $h^A_{(g)}$ is non-increasing by construction; so each $f_A$ is concave.

Maximising a separable sum of concave functions subject to a cardinality constraint is exactly the problem solved by the greedy rule of repeatedly taking the largest remaining marginal gain, which is optimal for separable concave resource allocation.
Since each marginal gain is computable in constant time after sorting each group, the total cost is $O(|\mathcal{C}'|\log|\mathcal{C}'|)$ for $|\mathcal{C}'|$ candidates.
\end{proof}

\alon{Done: this replaces the previous NP-hardness claim, whose reduction was invalid --- each candidate $(p,A)$ names a single attribute, so $k$ candidates cover at most $k$ attributes and the construction could not encode Maximum Coverage. The stronger exactness result above was verified by exhaustive search against greedy on $400$ random instances, matching in all of them.}

The theorem is about the explanation stage alone, with its candidate set and its scores held fixed, and it is worth being explicit about what it therefore does \emph{not} say.
It does not say that the two-stage framework maximises $F_w$ over all of $\mathcal{C}$: the view stage filters by $\Use$, which $F_w$ does not see, so a candidate discarded there is never compared under $F_w$ at all, and Lemma~\ref{lem:filtering} is what bounds the cost of that.
Nor does it apply across the rounds of Algorithm~\ref{alg:topk} as such, since $\gamma_i$ changes from round to round, so the scores compared in different rounds are not those of a single $F_w$.
What the theorem does give is the reference point used in Section~\ref{subsec:quality}: with the candidate set and $\gamma$ held fixed, greedy over $F_w$ \emph{is} the optimum, which is why the non-private greedy baseline there attains the exhaustive optimum on every query pair tested --- not slack in the $(1-1/e)$ bound, but the guarantee of the theorem --- and why the gap the private algorithm exhibits against it is attributable to the privacy mechanism.

\subsection{Private Query-Difference Explanation Algorithm}
\label{subsec:private-algorithm}

\begin{algorithm}[t]
\caption{Private-Queries-Diff-Top-$k$-Explanation}
\label{alg:main}
\begin{algorithmic}[1]
\Require Private dataset $D$; queries $Q_1,Q_2$; views $\mathcal{P}$; explanation attributes $\Attrs\subseteq\Attr(D)$; view budget $\tau_{\max}$; explanations $k$; smoothing constants $\gamma_{\text{ref}}$ and $\gamma_R$; diversity $\lambda$; privacy budget $\epsilon$
\Ensure $k$ differentially private explanation histograms
\State $\epsilon_P,\epsilon_S,\epsilon_H \gets 0.25\epsilon,\ 0.45\epsilon,\ 0.30\epsilon$
\State \textbf{/* Private view selection */}
\State $u_{use}(p)\gets \Use(p,D,Q_1,Q_2)$
\State $\widehat{\mathcal{P}}\gets \textsc{OneShotTop-}\tau(\mathcal{P},u_{use},\epsilon_P,\Delta{=}\tfrac{9}{4|D|})$
\Comment{$\epsilon_P$-DP}
\Statex \textbf{(Post-processing)} Candidate construction incurs no privacy loss.
\State $\mathcal{E}\gets \textsc{Find-Top-$k$-Expl}(D,Q_1,Q_2,\widehat{\mathcal{P}},$
\Statex \hspace{\algorithmicindent}$\Attrs,k,\gamma_{\text{ref}},\epsilon_S,\epsilon_H)$
\State \Return $\mathcal{E}$
\end{algorithmic}
\end{algorithm}

\begin{algorithm}[t]
\caption{Find-Top-$k$-Explanations}
\label{alg:topk}
\begin{algorithmic}[1]
\Require $D$; $Q_1,Q_2$; selected views $\widehat{\mathcal{P}}$; attributes $\Attrs$; $k$; $\gamma_{\text{ref}}$; $\tau_{\max}$; diversity $\lambda$; budgets $\epsilon_S,\epsilon_H$
\State $Counts\gets\emptyset$;\ \ $\mathcal{E}\gets\emptyset$;\ \ $\epsilon_i \gets \epsilon_S/k$
\For{$i\gets 1$ to $k$}
  \State $\tau_i \gets$ \Call{AdaptiveTau}{$\epsilon_{\text{used}},\epsilon_S,\tau_{\max}$} \Comment{public accountant}
  \State $\gamma_i \gets \gamma_{\text{ref}}\cdot \epsilon_{\text{ref}}/\epsilon_i$ \Comment{public; no privacy cost}
  \State $C_i \gets \{(p,A): p\in\widehat{\mathcal{P}}[1..\tau_i],\ A\in \Attrs\}\setminus \mathcal{E}$
  \State $u_i(p,A)\gets \HD_{\gamma_i}(D;p,A,Q_1,Q_2)\cdot \lambda^{Counts[A]}$
  \State $c_i=(p_i,A_i)\gets \ExpMech\!\left(C_i,u_i,\epsilon_i,\tfrac{1}{1+\gamma_i}\right)$
  \State $H^1_i\gets \hist{Q_1(p_i(D))}{A_i}$;\ \ $H^2_i\gets \hist{Q_2(p_i(D))}{A_i}$
  \State $\hat H^j_i\gets H^j_i+\Lap(2k/\epsilon_H)$ \textbf{for} $j=1,2$
  \State $\mathcal{E}\gets\mathcal{E}\cup\{(p_i,A_i,\hat H^1_i,\hat H^2_i)\}$;\ \ $Counts[A_i]{+}{+}$
\EndFor
\State \Return $\mathcal{E}$
\end{algorithmic}
\end{algorithm}

The $k$ rounds of Algorithm~\ref{alg:topk} implement the one-shot top-$k$ mechanism of Definition~\ref{def:oneshot-topk}, applied with the smoothed sensitivity of Proposition~\ref{prop:hd-sensitivity}.
Note that the rounds are written out explicitly here rather than as a single call, because the diversity factor makes the utility of round $i$ depend on the attributes selected in rounds $1,\dots,i-1$; the one-shot form applies unchanged to the view-selection stage, where the utility is fixed in advance.

\subsection{Differential Privacy Analysis}
\label{subsec:dp-proof}

\begin{proposition}[Sensitivity of the released histograms]
\label{prop:histogram-release-sensitivity}
Fix a candidate $(p,A)$. The pair $\big(\hist{Q_1(\pD)}{A},\hist{Q_2(\pD)}{A}\big)$ has $L_1$-sensitivity at most $2$, and their difference $\Delta H_{p,A}$ has $L_1$-sensitivity at most $1$.
\end{proposition}

\begin{proof}
One tuple joins at most one bin of $A$, and it may be selected by both queries, by one of them, or by neither.
It therefore changes each of the two histograms in at most one coordinate and by at most $1$, which gives $2$ for the pair.
For the difference, the affected coordinate changes by $d_1-d_2$ with $d_1,d_2\in\{0,1\}$, so by at most $1$, and all other coordinates are unchanged.
\end{proof}

Releasing the pair therefore calls for $\Lap(2k/\epsilon_H)$, as in Algorithm~\ref{alg:topk}; the variant that releases only the difference uses $\Lap(k/\epsilon_H)$ for the same budget.

\begin{proposition}[Differential Privacy of Algorithm~\ref{alg:topk}]
\label{thm:alg2-dp}
Algorithm~\ref{alg:topk} satisfies $(\epsilon_S+\epsilon_H)$-differential privacy.
\end{proposition}

\begin{proof}
Each round invokes the exponential mechanism with budget $\epsilon_S/k$ on a utility of sensitivity $1/(1+\gamma_i)$ (Proposition~\ref{prop:selection-utility-sensitivity}), which is $(\epsilon_S/k)$-DP, and the Laplace mechanism with budget $\epsilon_H/k$ on a pair of histograms of $L_1$-sensitivity $2$ (Proposition~\ref{prop:histogram-release-sensitivity}), which is $(\epsilon_H/k)$-DP.
The quantities $\tau_i$ and $\gamma_i$ are functions of the public accountant only and are therefore not privacy-relevant.
The rounds are adaptive, since $u_i$ depends on previously selected attributes, but adaptive sequential composition applies \cite{DBLP:series/synthesis/2016Li}, giving $k(\epsilon_S/k+\epsilon_H/k)=\epsilon_S+\epsilon_H$.
\end{proof}

\begin{proposition}[Differential Privacy of Algorithm~\ref{alg:main}]
\label{thm:alg1-dp}
Algorithm~\ref{alg:main} satisfies $\epsilon$-differential privacy.
\end{proposition}

\begin{proof}
View selection is $\epsilon_P$-DP by Proposition~\ref{prop:usefulness-sensitivity}, which supplies the sensitivity $9/(4|D|)$ that Algorithm~\ref{alg:main} passes to the mechanism, together with Theorem~\ref{thm:oneshot-topk}.
Candidate construction is post-processing.
Algorithm~\ref{alg:topk} is $(\epsilon_S+\epsilon_H)$-DP by Proposition~\ref{thm:alg2-dp}.
Sequential composition gives $\epsilon_P+\epsilon_S+\epsilon_H=\epsilon$.
\end{proof}

\begin{corollary}[End-to-End Differential Privacy]
\label{cor:end-to-end-dp}
The overall explanation pipeline satisfies $\epsilon$-differential privacy.
\end{corollary}

\subsection{Complexity Analysis}
\label{subsec:complexity-analysis}

Let $|\mathcal{P}|$ be the number of candidate views, $a=|\Attrs|$ the number of explanation attributes, $b$ the maximum number of bins per attribute, and $N=|D|$.
(We write $a$ rather than $m$ here because $m$ already denotes the view size $|p(D)|$ in Section~\ref{sec:explanation-quality}.)
Computing $\Use$ for every view costs $O(|\mathcal{P}|\,N)$, and the one-shot top-$\tau$ mechanism adds $O(\tau|\mathcal{P}|)$.
Each of the $k$ rounds scores at most $\tau a$ candidates and releases one histogram, giving $O(k\,\tau\,a\,b)$ once the per-candidate bin vectors are materialised, plus $O(\tau a\, N)$ to materialise them.
The total is
\[
O\!\left(|\mathcal{P}|\,N \;+\; \tau a\, N \;+\; k\,\tau\,a\,b\right).
\]
The two $O(N)$ terms are a one-off preprocessing cost independent of $\epsilon$ and $k$; only the third is paid per round.
Because $\HD_\gamma$ depends on $\gamma$ solely through the denominator, caching the per-bin numerator and denominator vectors makes the score computable for \emph{any} $\gamma$ without revisiting the data --- which is what keeps the adaptive $\gamma_i$ of Equation~\eqref{eq:adaptive-params} affordable, and is reflected in the measured separation of build and per-round cost in Section~\ref{subsec:scalability}.
